Let $d$ be the maximum distance between any two vertices of the polygon,
and let $A$ and $B$ be two vertices at distance $d$.
The vertices of the polygon, and therefore also the polygon itself,
are contained in the strip delimited by the two lines perpendicular to $\overline{AB}$ that pass, respectively, through $A$ and through $B$.
This is because any point outside of that region has distance to $A$ or to $B$ greater than $d$.

Let $C$ be a vertex of the polygon at maximal distance to the line $AB$,
and let $r$ be that distance.
The vertices of the polygon, and therefore also the polygon itself,
are contained in the strip delimited by the two lines parallel to $\overline{AB}$ at distance $r$ from the line $AB$.
This is because any point outside of that region has distance to $AB$ greater than $r$.

Therefore, the polygon is contained in the intersection of the two strips,
which is a rectangle of side lengths $d$ and $2 r$.

The triangle $ABC$ is contained in the polygon (by convexity),
and has area $\mfrac{1}{2} r d$.
The polygon has area less than or equal to $2 r d$
(the area of the rectangle containing it).
Hence, the triangle's area is at least one fourth of the polygon's area.
